\begin{definition}[Trace length]
\label{def:trace-length}
For a closed triangle $T$ and a rectifiable set $E$, define
\[
 L_E(T):=\mathcal H^1(T\cap E).
\]
If the seven open triangles cover $E$, their closures satisfy
\begin{equation}
 \mathcal H^1(E)
 \le L_E(T_C)+\sum_{i=0}^5L_E(T_i).
 \label{eq:length-subadditivity}
\end{equation}
\end{definition}

The two targets are the perimeter $\partial H$ and the full skeleton $S$.
Their lengths are $6$ and $12$, respectively.

\begin{figure}[H]
\centering
\begin{tikzpicture}[scale=.84]
  \foreach \x/\title in {0/$\partial H$: length $6$,4/$D$: length $6$,8/$S=\partial H\cup D$: length $12$}{
    \begin{scope}[xshift=\x cm]
      \coordinate (a0) at (1,0);
      \coordinate (a1) at (.5,.866);
      \coordinate (a2) at (-.5,.866);
      \coordinate (a3) at (-1,0);
      \coordinate (a4) at (-.5,-.866);
      \coordinate (a5) at (.5,-.866);
      \coordinate (o) at (0,0);
      \draw[hex] (a0)--(a1)--(a2)--(a3)--(a4)--(a5)--cycle;
      \node[panellabel,align=center] at (0,-1.38) {\title};
    \end{scope}
  }
  \draw[demand] (1,0)--(.5,.866)--(-.5,.866)--(-1,0)--(-.5,-.866)--(.5,-.866)--cycle;
  \foreach \p in {(4,0),(4.5,.866),(3.5,.866),(3,0),(3.5,-.866),(4.5,-.866)}{
    \draw[radialdemand] (4,0)--\p;
  }
  \draw[demand] (9,0)--(8.5,.866)--(7.5,.866)--(7,0)--(7.5,-.866)--(8.5,-.866)--cycle;
  \foreach \p in {(9,0),(8.5,.866),(7.5,.866),(7,0),(7.5,-.866),(8.5,-.866)}{
    \draw[radialdemand] (8,0)--\p;
  }
  \node[figlabel,align=center,text=DemandRed] at (8,-1.93)
    {the full-skeleton estimate is used only under its stated branch hypotheses};
\end{tikzpicture}

\caption{The perimeter and the full skeleton.  The skeleton is used only
under the hypotheses of the corresponding trace bound.}
\label{fig:strategy1-targets}
\end{figure}

\begin{theorem}[Boundary trace table]
\label{thm:boundary-trace-table}
For the closures of the original open triangles, the following bounds hold:
\begin{equation}
\begin{array}{c|c|c}
\text{triangle}&\text{hypothesis}&L_{\partial H}\\ \hline
T_C&N_E(T_C)=0&0\\
T_C&N_E(T_C)=1&\le\dfrac{\sqrt3}{2}-\dfrac34\\[3pt]
T_C&N_E(T_C)=2&<\dfrac12\\[3pt]
T_i&(o(T_i),n(T_i))\in\{(1,0),(2,0)\},\ A_i+B_i\le1&\le1\\
T_i&(o(T_i),n(T_i))\in\{(1,0),(2,0)\},\ A_i+B_i>1&\le\dfrac2{\sqrt3}\\[3pt]
T_i&(o(T_i),n(T_i))\in\{(1,1),(1,2)\}&<\dfrac12\\
T_i&(o(T_i),n(T_i))=(2,1)&<1.
\end{array}
\label{eq:boundary-trace-table}
\end{equation}
\end{theorem}

\begin{proof}
This is the exact output of the boundary-trace optimization in
\zcref{lem:app-boundary-trace-table-calculation}.
\end{proof}

\begin{lemma}[A positive-support adjacent exceptional triangle is forced]
\label{lem:positive-support-rescuer}
Suppose $N_E(T_C)\in\{1,2\}$ and, after symmetry,
\[
 T_C\cap\{M_0,\ldots,M_5\}=\{M_0\}.
\]
If $N_+\ge2$, then there are a supercritical index
$s\in\mathbb Z/6\mathbb Z\setminus\{0\}$ and
$j\in\{s-1,s+1\}$ such that
\[
 M_s\in U_j,
 \qquad
 \mathcal H^1(T_j\cap r_s)>0.
\]
\end{lemma}

\begin{proof}
Choose a supercritical $s\ne0$.  By Lemma~\ref{lem:self-midpoint},
$M_s\notin T_s$, and the displayed midpoint hypothesis gives
$M_s\notin T_C$.  Hence the original open cover puts $M_s$ in some $U_j$.
The diameter-one restriction gives $j\in\{s-1,s+1\}$, and openness gives
$\mathcal H^1(T_j\cap r_s)>0$.
\end{proof}

\begin{theorem}[Direct CE1/CE2 skeleton count]
\label{thm:common-skeleton-count}
If
\[
 N_E(T_C)\in\{1,2\},
 \qquad
 T_C\cap\{M_0,\ldots,M_5\}=\{M_0\},
 \qquad
 N_++N_{\mathrm{sp}}\ge3,
\]
then the seven triangles do not cover $S$.
\end{theorem}

\begin{proof}
The four optimized skeleton caps are
\zcref{lem:center-skeleton-cap},
\zcref{lem:supercritical-skeleton-cap},
\zcref{lem:positive-support-skeleton-cap}, and
\zcref{lem:no-support-skeleton-cap}.  The classes counted by $N_+$ and
$N_{\mathrm{sp}}$ are disjoint.  Therefore
\[
 L_S(T_C)+\sum_{i=0}^5L_S(T_i)
 <\frac32+\frac32(N_++N_{\mathrm{sp}})
   +2(6-N_+-N_{\mathrm{sp}})
 \le12=\mathcal H^1(S),
\]
contradicting \eqref{eq:length-subadditivity}.
\end{proof}

\begin{proposition}[Vd2 adjacent-midpoint branch]
\label{prop:ce2-vd2-midpoint-length}
Assume
\[
 N_E(T_C)=2,
 \qquad N_+=1,
 \qquad
 \{j:(o(T_j),n(T_j))\in\{(1,1),(1,2)\}\}=\{i\},
\]
\[
 (o(T_i),n(T_i))=(1,2),
 \qquad
 M_{i-1}\in T_i\ \lor\ M_{i+1}\in T_i,
\]
and
\[
 \forall j\ne i\qquad
 (o(T_j),n(T_j))\in\{(1,0),(2,0)\}.
\]
Then the boundary cannot be covered.
\end{proposition}

\begin{proof}
The exact adjacent-midpoint cap is
\zcref{lem:vd2-neighbor-midpoint-cap}; its substitution in the perimeter
budget is item~(5) of \zcref{cor:signed-budget-branches}.
\end{proof}

\begin{proposition}[Branches closed by trace length]
\label{prop:length-branches}
The trace-length mechanism excludes exactly the following routing entries
and hybrid subcase:
\begin{enumerate}
\item $N_{\mathrm{gap}}=0$ and $N_+=0$;
\item $N_{\mathrm{gap}}=0$, $N_+=1$, and
      $\#\{i:T_i\text{ is }\Vdone\text{ or }\Vdtwo\}\ge1$;
\item $N_E(T_C)\in\{1,2\}$ and $N_++N_{\mathrm{sp}}\ge3$;
\item $N_E(T_C)\in\{1,2\}$, $N_+=0$, and
      $\#\{i:T_i\text{ is }\Vdone\text{ or }\Vdtwo\}\ge1$;
\item
      \[
       \begin{gathered}
       N_E(T_C)=1,\qquad N_+=1,\\
       \#\{i:T_i\text{ is }\Vdone\text{ or }\Vdtwo\}=1,\qquad
       \#\{i:T_i\text{ is }\Tlike\}=0;
       \end{gathered}
      \]
\item the subcase in Proposition~\ref{prop:ce2-vd2-midpoint-length}.
\end{enumerate}
\end{proposition}

\begin{proof}
The complete exact substitutions, including every strict inequality, are
given in \zcref{lem:app-length-branch-calculation}.  Item~(3) is also the
direct consequence of Theorem~\ref{thm:common-skeleton-count}, and item~(6)
is Proposition~\ref{prop:ce2-vd2-midpoint-length}.
\end{proof}
